\documentclass[11pt]{article}

\usepackage[margin=1in]{geometry}
\usepackage[T1]{fontenc}
\usepackage[utf8]{inputenc}
\usepackage{lmodern}
\usepackage{microtype}
\usepackage{amsmath,amssymb,amsthm,amscd,mathtools}
\usepackage{booktabs}
\usepackage{enumitem}
\usepackage{tabularx}
\usepackage[hidelinks]{hyperref}
\usepackage[nameinlink,capitalise]{cleveref}

\newtheorem{theorem}{Theorem}[section]
\newtheorem{proposition}[theorem]{Proposition}
\newtheorem{lemma}[theorem]{Lemma}
\newtheorem{corollary}[theorem]{Corollary}
\theoremstyle{definition}
\newtheorem{definition}[theorem]{Definition}
\newtheorem{assumption}[theorem]{Assumption}
\newtheorem{example}[theorem]{Example}
\theoremstyle{remark}
\newtheorem{remark}[theorem]{Remark}

\newcommand{\dd}{\mathrm d}
\newcommand{\cA}{\mathcal A}
\newcommand{\cC}{\mathcal C}
\newcommand{\cD}{\mathcal D}
\newcommand{\cE}{\mathcal E}
\newcommand{\cF}{\mathcal F}
\newcommand{\cG}{\mathcal G}
\newcommand{\cH}{\mathcal H}
\newcommand{\cI}{\mathcal I}
\newcommand{\cK}{\mathcal K}
\newcommand{\cM}{\mathcal M}
\newcommand{\cN}{\mathcal N}
\newcommand{\cO}{\mathcal O}
\newcommand{\cP}{\mathcal P}
\newcommand{\cR}{\mathcal R}
\newcommand{\cS}{\mathcal S}
\newcommand{\cU}{\mathcal U}
\newcommand{\cV}{\mathcal V}
\newcommand{\cX}{\mathcal X}
\newcommand{\cY}{\mathcal Y}
\newcommand{\MTT}{Modal Triplet Theory}
\newcommand{\Ker}{\operatorname{Ker}}
\newcommand{\Ran}{\operatorname{Ran}}
\newcommand{\Vol}{\operatorname{Vol}}
\newcommand{\Tr}{\operatorname{Tr}}
\newcommand{\Id}{\operatorname{Id}}

% Keep multi-paragraph managed evidence rows on separate list lines.
\let\MTToriginalitem\item
\renewcommand{\item}{\par\MTToriginalitem}

\title{\vspace{-1em}
Modal Triplet Theory and Kaluza--Klein Reduction:\\[0.25em]
\large A Conditional \(4+6\) Spectral and Truncation Contract}
\author{Peter Nero}
\date{Version 2, July 2026}

\begin{document}
\maketitle

\begin{abstract}
This paper asks what Modal Triplet Theory (MTT) presently establishes about
Kaluza--Klein reduction.  A Kaluza--Klein construction begins with a supplied
higher-dimensional geometry, action, field content, operator domains, mode
normalizations, and reduction ansatz.  Compactness then gives discrete
internal spectra, but it does not select the internal manifold, its scale,
the higher-dimensional action, or a nonlinearly consistent truncation.  We
package the required inputs as a typed Kaluza--Klein record and define a
partial map from an upper MTT state to that record.  The main realization
theorem is conditional: if one selected MTT state emits every row of the
record, the lower consistency and error conditions hold, and MTT evaluation
factors through the standard reduction, then MTT realizes that reduced sector
at the declared cutoff and order.  For a fixed compact fiber we derive the
usual tower masses from the eigenvalues of the appropriate internal
Laplace-, Dirac-, or Lichnerowicz-type operator.  We also distinguish an
exact consistent truncation from a spectral projection and give a
Lyapunov--Schmidt residual certificate for controlled approximate reduction.
The canonical MTT specialization is a ten-dimensional bundle over a
four-dimensional Lorentzian base with compact six-dimensional Riemannian
fiber.  Its constant-time \(9\to3\) description is the spatial slice of the
same \(10\to4\) fibration, not an independent derivation.  The local
\(1+3\times3=4+6\) component identity and the selected finite shared-line
carrier motivate this geometry, but they do not yet construct the physical
six-manifold or preserve its connection and Hessian.  Consequently the
paper establishes a reusable and testable reduction contract, not a
first-principles selection of extra dimensions, Kaluza--Klein radii, gauge
groups, particle masses, or Standard Model parameters.
\end{abstract}

\section*{Revision note: Version 2}
\addcontentsline{toc}{section}{Revision note: Version 2}
\begin{description}[leftmargin=2.5cm,style=nextline]
\item[Supersedes]
Version 2 supersedes Version 1 and its claims that MTT had already selected
the Kaluza--Klein internal space, radius, zero-mode spectrum, gauge
couplings, family number, and a finite fixed-point system equivalent to the
full higher-dimensional field equations.

\item[Reason]
The former paper combined valid Kaluza--Klein formulas with unproved MTT
source claims.  It treated a finite left-invariant ansatz as automatically
complete for a nonlinear PDE, used cohomological completeness as if it were
solution-space completeness, and inferred controlled truncation from a gap
ratio without bounding the omitted equations.  It also described the
constant-time \(9\to3\) slice as an independently proved compactification
equivalence and imported obsolete Iwasawa and Lens--Nil selection claims.

\item[Resolution]
This revision defines the full lower record, separates kinematic mode
expansion from dynamical truncation, proves a conditional MTT realization
theorem, and gives an explicit omitted-mode residual certificate.  It uses
the canonical \(4+6\) geometry and states precisely what the
\(1+3\times3=4+6\) identity does and does not imply.  Standard
Kaluza--Klein formulas retain their literature ownership.

\item[Retained content]
Compact internal operators still organize four-dimensional towers; internal
isometries can supply gauge fields; normalized overlap integrals determine
effective couplings; and finite invariant ansatzes can exactly encode the
equations restricted to those ansatzes.  These are valuable lower-level
tools once their inputs and domains are supplied.

\item[Remaining boundary]
The current MTT research ledger leaves the continuum
geometry-to-operator naturality problem open.  Its exit requires a
same-source commuting map from the world-in-world/strain carrier to the
physical \(q=79\) vertical Hermitian--Yang--Mills complex that preserves the
connection, covariant derivatives, and Hessian.  Without that map, the
finite carrier does not select the physical Kaluza--Klein fiber or its
operators.
\end{description}

\tableofcontents

\section{The corrected question}
\label{sec:question}

Kaluza--Klein theory is a method of reduction.  One starts with a field
theory on a space of dimension \(D=4+d\), resolves its fields into internal
modes, and asks which lower-dimensional fields and interactions survive at
the scale of interest.  Kaluza's five-dimensional metric ansatz and Klein's
compact circle gave the original example; modern supergravity and string
compactifications use the same architecture with richer fields and internal
geometries
\cite{Kaluza1921,Klein1926,DuffNilssonPope1986,Witten1981}.

There are two logically different questions:
\begin{enumerate}[leftmargin=2em]
\item \emph{Reduction question.}  Given the higher-dimensional theory and
      compactification data, what lower-dimensional theory follows?
\item \emph{Source question.}  Why are that theory, internal geometry,
      radius, bundle, vacuum, and truncation selected?
\end{enumerate}
Kaluza--Klein analysis answers the first question.  It does not by itself
answer the second.  The corrected MTT question is therefore:
\begin{quote}
Does one selected upper MTT state emit a complete higher-dimensional record,
and does MTT evaluation factor through a mathematically controlled
Kaluza--Klein reduction of that same record?
\end{quote}

This formulation makes the burden of proof visible.  A metric with six
unlabelled internal coordinates is not yet a compactification.  Nor is an
orthogonal projector a consistent low-energy theory.  The action, gauge
identifications, operator domains, mode normalization, nonlinear ansatz,
source terms, and approximation error all matter.

\subsection{Result and non-result}

The positive result of this paper is a conditional realization contract.  It
states exactly what MTT must provide and what follows once it does.  The
spectral and truncation statements below are ordinary mathematical
consequences of a fixed lower record.  They are included because they expose
the interface that an MTT source theorem must satisfy.

The paper does not prove that MTT already selects a unique compact
six-manifold, radius, higher-dimensional action, gauge group, chiral
spectrum, or observed mass.  In particular, a dimension count is not a
geometric construction, and agreement with a profile used to choose a
branch is not a held-out prediction.

\section{The lower Kaluza--Klein object}
\label{sec:record}

\subsection{Geometry}

Let \(M_4\) be a time-oriented Lorentzian four-manifold and let
\[
  \pi:Y_D\longrightarrow M_4,\qquad D=4+d,
\]
be a smooth fiber bundle with compact \(d\)-dimensional fibers
\(X_x=\pi^{-1}(x)\).  A horizontal distribution
\(\cH\subset TY_D\) gives
\[
  TY_D=\cH\oplus\cV,\qquad \cV=\Ker\dd\pi .
\]
The higher-dimensional metric is Lorentzian on the horizontal directions
and Riemannian on the vertical directions.  In a local trivialization it may
be written schematically as
\begin{equation}
\label{eq:kkmetric}
\dd s_D^2
 =g_{\mu\nu}(x)\,\dd x^\mu\dd x^\nu
 +h_{mn}(x,y)
   \bigl(\dd y^m+\cA^m{}_\mu(x,y)\dd x^\mu\bigr)
   \bigl(\dd y^n+\cA^n{}_\nu(x,y)\dd x^\nu\bigr).
\end{equation}
Equation \eqref{eq:kkmetric} is a local decomposition of supplied metric
data.  It neither proves that \(Y_D\) exists globally nor selects
\((X,h,\cH)\).

Warping, boundaries, orbifold strata, monodromy, and varying fibers require
additional data.  To keep the spectral theorems transparent, the principal
calculation below uses a closed compact fiber and a product or
adiabatically controlled background.  More general cases require the
corresponding family, boundary, or singular elliptic theory.

\subsection{Fields, action, and operators}

A reduction must specify the higher-dimensional fields and the action from
which their equations follow.  We denote by \(\cF_D\) the collection of
metrics, connections, spinors, differential forms, scalars, source data, and
any gauge-fixing or constraint variables.  Let
\[
 S_D:\cD(S_D)\subset\cF_D\longrightarrow\mathbb R
\]
be the action at a declared derivative and quantum order, with equations
\(\cE_D(\Phi)=0\).  Gauge transformations, diffeomorphisms, boundary
conditions, and normalizations are part of the record rather than
ellipsis.

Each field type has its own vertical operator.  Scalar modes may be
organized by a Laplace-type operator, spinors by a twisted Dirac operator,
one-forms by a gauge-fixed Hodge operator, and metric fluctuations by a
Lichnerowicz-type operator.  Background flux, curvature, torsion, and
potentials may add lower-order terms.  A single formula
\(-\Delta_XY_n=\lambda_nY_n\) is therefore illustrative, not universal.

\begin{definition}[Complete Kaluza--Klein record]
\label{def:record}
A \emph{complete Kaluza--Klein record} at scope
\(\sigma=(D,N,\ell,\partial,\mathrm{obs})\) is
\[
\cK_\sigma=
\bigl(
Y_D,\pi,g_D,\cH;\,
\cF_D,S_D,\cG_D;\,
\{L_\alpha,\cD(L_\alpha)\}_\alpha;\,
\{u_{\alpha n}\};\,
P_N,\iota_N,\cE_4;\,
\cC_N,\cN,\cO
\bigr),
\]
where:
\begin{enumerate}[label=(\roman*),leftmargin=2em]
\item \((Y_D,\pi,g_D,\cH)\) is the global geometric and connection data;
\item \((\cF_D,S_D,\cG_D)\) gives fields, action, symmetries, constraints,
      sources, and boundary conditions;
\item every \(L_\alpha\) is the appropriate closed vertical operator with a
      specified domain, measure, adjoint convention, and normalization;
\item \(\{u_{\alpha n}\}\) is a complete normalized mode system whenever
      such a system is invoked;
\item \(P_N\) is the retained spectral or representation projector and
      \(\iota_N\) is the full nonlinear uplift ansatz, not merely a linear
      inclusion;
\item \(\cE_4\) is the claimed lower equation or action;
\item \(\cC_N\) records exact consistency or an explicit residual and error
      certificate;
\item \(\cN\) contains coupling, scale, frame, and measure normalizations;
      and
\item \(\cO\) declares the observables and comparison map at which the
      reduction is asserted.
\end{enumerate}
\end{definition}

The record is deliberately longer than a metric ansatz.  Most false
compactification arguments omit precisely the rows that decide whether the
calculation is physical: the nonlinear uplift, normalization, domain, and
error rows.

\section{The MTT realization contract}
\label{sec:contract}

Let \(\cU_{\mathrm{MTT}}\) be the domain of admissible upper MTT states.  A
candidate source map is a partial typed map
\[
  \cR_{\mathrm{KK}}:
  \cU_{\mathrm{MTT}}\dashrightarrow \mathfrak K_\sigma ,
\]
where \(\mathfrak K_\sigma\) is the class of records in
Definition~\ref{def:record}.  The dashed arrow matters: not every abstract coherent
state contains spacetime, action, or reduction data.

Let
\[
 \operatorname{Red}_{\mathrm{KK}}:\mathfrak K_\sigma
 \dashrightarrow \mathfrak L_\sigma
\]
be the standard lower-dimensional reduction functional and let
\(\operatorname{Ev}_{\mathrm{MTT}}\) and
\(\operatorname{Ev}_4\) be the upper and lower observable evaluations on
their declared common domain.

\begin{definition}[Complete MTT--KK realization]
\label{def:realization}
A selected state \(u_\ast\in\cU_{\mathrm{MTT}}\) realizes a
Kaluza--Klein sector at scope \(\sigma\) when:
\begin{enumerate}[label=(\roman*),leftmargin=2em]
\item \(\cR_{\mathrm{KK}}(u_\ast)\) is a complete record;
\item every global, gauge, analytic, and source condition required by that
      record is satisfied;
\item the truncation certificate \(\cC_N\) is exact or gives a stated error
      bound on the requested domain; and
\item the observable diagram commutes,
\[
\begin{CD}
u_\ast @>{\cR_{\mathrm{KK}}}>>
       \cK_\sigma @>{\operatorname{Red}_{\mathrm{KK}}}>>
       \mathfrak L_\sigma\\
@V{\operatorname{Ev}_{\mathrm{MTT}}}VV
&& @VV{\operatorname{Ev}_4}V\\
\cO_\sigma @= \cO_\sigma .
\end{CD}
\]
\end{enumerate}
\end{definition}

Equivalently, the commuting condition may be read algebraically as
\[
 \operatorname{Ev}_{\mathrm{MTT}}(u_\ast)
 =
 \operatorname{Ev}_4\!\left(
 \operatorname{Red}_{\mathrm{KK}}
 (\cR_{\mathrm{KK}}(u_\ast))\right).
\]

\begin{theorem}[Conditional MTT--Kaluza--Klein realization]
\label{thm:conditional}
Suppose a selected MTT state \(u_\ast\) satisfies
Definition~\ref{def:realization}.  Then MTT realizes the lower
Kaluza--Klein sector
\[
 \operatorname{Red}_{\mathrm{KK}}
 \bigl(\cR_{\mathrm{KK}}(u_\ast)\bigr)
\]
at scope \(\sigma\).  If \(\cC_N\) is exact, the assertion is exact on the
declared ansatz.  If \(\cC_N\) supplies error \(\epsilon_N\), every observable
in \(\cO_\sigma\) is realized only to the propagated stated error.
\end{theorem}

\begin{proof}
The source map supplies one lower record rather than a list of mutually
unrelated compatible objects.  The standard reduction functional is
therefore defined on that record.  The consistency certificate makes the
reduction exact or controlled at the declared scope, while commutation of
the evaluation diagram identifies the upper and lower observable values.
No conclusion is asserted outside the domains, order, cutoff, sources, or
observables recorded in \(\sigma\).
\end{proof}

\begin{remark}[Why the theorem is substantive but conditional]
The theorem is not the empty statement that equal models are equal.  It
fixes the exact source, covariance, normalization, truncation, and evaluation
obligations whose omission previously allowed a representation to be
remembered as a derivation.  Its hypotheses are nevertheless not yet all
discharged by current MTT geometry.
\end{remark}

\section{Canonical \(4+6\) geometry}
\label{sec:fourplus}

\subsection{The physical specialization}

The canonical physical specialization used in current MTT is
\[
  \pi:M_{10}\longrightarrow M_4,
  \qquad \dim M_4=4,\qquad \dim X_x=6,
\]
with a globally hyperbolic Lorentzian base in the physical completion and
compact Riemannian fibers.  Positive elliptic modal operators act
vertically.  They organize internal states; they do not add causal time
directions \cite{NeroFoundation}.

This \(4+6\) architecture is a declared physical realization of the
dimension-neutral MTT Hilbert-bundle formalism.  It is not derived merely
from the existence of three modal labels.  If a selected physical
six-manifold is eventually supplied, it must still come with the geometry,
connections, action, operators, and reduction data in
Definition~\ref{def:record}.

\subsection{What the \(3\times3\) field actually counts}

Let \(TP\) and \(TI\) be oriented Euclidean rank-three bundles.  A local
world-in-world comparison field
\[
 Q_{\mathrm{WW}}\in\Gamma(\operatorname{Hom}(TP,TI))
\]
has nine components after choosing frames.  Around a nonsingular background,
polar decomposition gives
\[
 \operatorname{Mat}(3,\mathbb R)
 =
 \mathfrak{so}(3)\oplus\operatorname{Sym}(3,\mathbb R),
 \qquad 9=3+6.
\]
Relative to an orthonormal flag,
\[
 \operatorname{Sym}(3,\mathbb R)
 =
 \mathbb RI_3\oplus\cD_0\oplus\cO,
 \qquad \dim(\mathbb RI_3,\cD_0,\cO)=(1,2,3).
\]
Thus the component identity
\[
 1+3\times3=(1+3)+(1+2+3)=4+6=10
\]
is exact once one ordering scalar is supplied.  It is not manifold-dimension
multiplication.  Nor does it prove that the four components are a Lorentzian
tangent space or that the six strain components globalize to the tangent
bundle of a compact physical fiber.

The identity is useful because it presents a candidate local interface:
orientation directions can be quotiented while the six strain directions
carry the \(1+2+3\) filtration.  The missing theorem is global.  It must
intertwine bundles, connections, covariant derivatives, measures, and
operators, not just match ranks.

\subsection{The constant-time \(9\to3\) picture}

Suppose the ten-dimensional fibration already has compatible splittings
\[
 M_{10}\simeq\mathbb R_t\times\Sigma_9,
 \qquad
 M_4\simeq\mathbb R_t\times M_3,
\]
and \(\pi\) preserves the time coordinate.  Restriction to a time slice then
gives
\[
 \pi_t:\Sigma_9\longrightarrow M_3
\]
with the same six-dimensional fiber.  Reattaching time recovers the original
\(10\to4\) fibration.

\begin{proposition}[Spatial-slice identity]
\label{prop:slice}
Under the compatible product and time-preservation hypotheses above, the
\(9\to3\) spatial projection and the \(10\to4\) fibration describe the same
fiberwise geometric data after restriction and reattachment of the supplied
time factor.
\end{proposition}

\begin{proof}
The restriction of \(\pi\) to
\(\{t\}\times\Sigma_9\) has base
\(\{t\}\times M_3\) and unchanged vertical bundle.  Conversely, taking the
product of \(\pi_t\) with \(\Id_{\mathbb R_t}\) reconstructs \(\pi\).  The
claim is therefore an identity between two descriptions of one already
supplied fibration.
\end{proof}

\begin{remark}
Proposition~\ref{prop:slice} is not a derivation of time, three-space, or six compact
directions.  If the metric is not product-like, the fibration depends on
time, or the horizontal distribution mixes causal and vertical directions,
additional hypotheses replace this elementary identity.
\end{remark}

\subsection{The shared circle}

Current MTT finite geometry contains one universal flat cyclic line whose
pullbacks agree across several selected finite carriers.  This is stronger
than noticing several isomorphic copies of \(U(1)\): it identifies common
connection and holonomy data at the proven finite level
\cite{NeroFoundation}.  It still does not make that line:
\begin{itemize}[leftmargin=2em,itemsep=0.2em]
\item \makebox[0.95\linewidth][l]{an extra seventh internal coordinate;}
\item \makebox[0.95\linewidth][l]{physical Lorentzian time;}
\item \makebox[0.95\linewidth][l]{the Kaluza--Klein circle of five-dimensional gravity; or}
\item \makebox[0.95\linewidth][l]{the M-theory reduction circle.}
\end{itemize}
Any of those identifications requires a source-preserving geometric
intertwiner with the relevant metric radius, spin structure, connection,
and action.  The corrected M-theory paper states the analogous circle
boundary for eleven-dimensional reduction \cite{NeroMTheory}.

\section{Spectral origin of Kaluza--Klein masses}
\label{sec:spectrum}

\subsection{A fixed compact fiber}

Let \(X\) be a closed compact Riemannian manifold and let \(E\to X\) be a
Hermitian vector bundle.  A formally self-adjoint elliptic operator
\[
 L_X:\cD(L_X)\subset L^2(X,E)\longrightarrow L^2(X,E)
\]
with its self-adjoint realization has compact resolvent.  Its spectrum is
discrete with finite multiplicities and an orthonormal eigenbasis
\(\{u_n\}\) \cite{Chavel1984,ReedSimonIV}.  The use of internal momentum and
twists to generate lower-dimensional masses is standard Kaluza--Klein
machinery \cite{ScherkSchwarz1979}.

Consider first a scalar on the unwarped product \(M_4\times X\) with
\[
 S[\Phi]=-\frac12\int_{M_4\times X}
 \left(
 |\dd_{M_4}\Phi|^2
 +\langle\Phi,L_X\Phi\rangle
 +m_D^2|\Phi|^2
 \right)\dd\mathrm{vol}_{4}\dd\mathrm{vol}_{X}.
\]
For
\[
 L_Xu_n=\lambda_nu_n,\qquad
 \int_X\overline{u_m}u_n\,\dd\mathrm{vol}_{X}=\delta_{mn},
\qquad
 \Phi(x,y)=\sum_n\phi_n(x)u_n(y),
\]
orthogonality diagonalizes the quadratic action.

\begin{lemma}[Fixed-background spectral mass dictionary]
\label{lem:masses}
Under the preceding hypotheses, the four-dimensional scalar modes have
\[
 m_n^2=m_D^2+\lambda_n.
\]
If \(h=R^2h_0\) and \(L_X=-\Delta_h\), then
\[
 \lambda_n(h)=R^{-2}\widehat\lambda_n(h_0),
 \qquad
 m_n^2=m_D^2+\frac{\widehat\lambda_n}{R^2}.
\]
The zero modes are \(\Ker L_X\); a positive first omitted eigenvalue gives a
quadratic spectral separation, not by itself a nonlinear truncation.
\end{lemma}

\begin{proof}
Insert the normalized eigenmode expansion into the quadratic action.
Self-adjointness and orthogonality remove cross terms, leaving one
four-dimensional quadratic action for each \(n\) with mass squared
\(m_D^2+\lambda_n\).  The inverse-square scaling follows from the metric
scaling of the scalar Laplacian.
\end{proof}

\subsection{Different fields use different operators}

For a twisted spinor, the internal Dirac eigenvalue enters the
four-dimensional mass matrix, with chirality and zero modes controlled by
the spin or spin\(^{c}\) bundle, twisting connection, and index.  For
differential forms, the relevant Hodge operator must be combined with gauge
constraints and harmonic representatives.  Metric fluctuations use a
gauge-fixed Lichnerowicz-type operator and can mix with form and scalar
fluctuations.  Background curvature, torsion, flux, or bundle curvature
changes the operator and may produce a matrix-valued eigenproblem.

It is therefore more accurate to write
\[
  M^2_{\alpha,mn}
  =
  M^2_{\alpha,\mathrm{bulk}}\delta_{mn}
  +
  \bigl\langle u_{\alpha m},
  L_{\alpha,X}u_{\alpha n}\bigr\rangle
\]
after all constraints and normalizations are fixed.  A numerical
Kaluza--Klein mass is a conclusion only after \(L_{\alpha,X}\), its domain,
the scale, and the background are selected independently of the target
mass.

\subsection{Uniform gaps are additional results}

Every fixed closed compact elliptic problem has discrete spectrum.  A
uniform lower bound on the first positive eigenvalue across a varying family
does not follow from compactness of each member.  Collapse, degeneration, or
an eigenvalue crossing zero can destroy such a bound.  The corrected
Calabi--Yau paper owns the corresponding fixed-background spectral
admissibility statement and its uniformity warning
\cite{NeroCalabiYau}.  The present paper uses that result rather than
duplicating it as an MTT theorem.

\section{Projection is not truncation}
\label{sec:truncation}

\subsection{The nonlinear consistency condition}

Let \(P_N\) project onto retained modes and let \(Q_N=\Id-P_N\).
Writing a field as \(P_N\Phi+Q_N\Phi\) is a kinematic decomposition.
Setting \(Q_N\Phi=0\) is dynamically valid only if the omitted equations
remain satisfied.

\begin{definition}[Exact consistent truncation]
\label{def:consistent}
An uplift \(\iota_N:\cF_4\to\cF_D\) is an exact consistent truncation when
there is a lower equation \(\cE_4\) such that
\[
 \cE_D(\iota_N\varphi)=0
 \quad\Longleftrightarrow\quad
 \cE_4(\varphi)=0
\]
on the declared domains.  Equivalently, every lower solution uplifts to a
higher-dimensional solution, and no omitted equation is sourced by the
retained configuration.
\end{definition}

Generic zero-mode truncations are not automatically consistent at nonlinear
order.  Products of retained harmonics can contain omitted harmonics.
Special sphere, group-manifold, and supergravity ansatzes may nevertheless
be consistent because symmetry and nonlinear field combinations enforce
cancellations \cite{CveticLuPope2000}.  The consistency belongs to the
specific uplift ansatz, not to the word ``compact.''

\subsection{A controlled residual certificate}

When exact consistency fails, one can still justify an effective reduction
by solving for the omitted modes and bounding their response.  The following
standard Lyapunov--Schmidt estimate makes the needed information explicit.

\begin{proposition}[Omitted-mode correction bound]
\label{prop:residual}
Let \(\cE_D:\cX\to\cY\) be \(C^1\) between Banach spaces with compatible
splittings \(P_N+Q_N=\Id\).  Fix a retained configuration
\(u_0=\iota_N\varphi\), set
\[
 r_N=Q_N\cE_D(u_0),\qquad
 A_N=Q_ND\cE_D(u_0)\big|_{Q_N\cX},
\]
and suppose \(A_N:Q_N\cX\to Q_N\cY\) is invertible with
\(\|A_N^{-1}\|\le\gamma_N^{-1}\).  On the ball
\(\|\eta\|\le\rho\), suppose
\[
\left\|
Q_N\!\left[
\cE_D(u_0+\eta)-\cE_D(u_0)-D\cE_D(u_0)\eta
\right]
-
Q_N\!\left[
\cE_D(u_0+\zeta)-\cE_D(u_0)-D\cE_D(u_0)\zeta
\right]
\right\|
\le L_\rho\|\eta-\zeta\|.
\]
If
\[
 \frac{L_\rho}{\gamma_N}<1,
 \qquad
 \frac{\|r_N\|}{\gamma_N}
 \le
 \left(1-\frac{L_\rho}{\gamma_N}\right)\rho ,
\]
then there is a unique \(\eta_\ast\in Q_N\cX\) in that ball such that
\[
 Q_N\cE_D(u_0+\eta_\ast)=0,
\qquad
 \|\eta_\ast\|
 \le
 \frac{\|r_N\|}
      {\gamma_N-L_\rho}.
\]
\end{proposition}

\begin{proof}
Write the omitted equation as the fixed-point problem
\[
\eta=
-A_N^{-1}r_N
-A_N^{-1}Q_N\!
\left[
\cE_D(u_0+\eta)-\cE_D(u_0)-D\cE_D(u_0)\eta
\right].
\]
The first inequality makes this map a contraction, and the second makes the
closed radius-\(\rho\) ball invariant.  Banach's fixed-point theorem gives
existence and uniqueness.  Taking norms in the fixed-point equation gives
the stated bound.
\end{proof}

\begin{remark}[What a gap does]
A spectral gap can help bound \(\|A_N^{-1}\|\), but the other quantities do
not disappear.  One must still compute the omitted residual \(r_N\), control
the nonlinear Lipschitz constant \(L_\rho\), and propagate the correction
into the retained equation and requested observables.  A dimensionless ratio
of two named scales is not an error estimate unless these links are proved.
\end{remark}

\section{Finite invariant ansatzes and the FCC}
\label{sec:fcc}

Version 1 introduced a fixed-point compactification condition (FCC): a
finite algebraic system obtained by expanding a higher-dimensional
background in a left-invariant basis.  That construction has a correct but
narrow scope.

Let \(\cA_{\mathrm{inv}}=\{u(a):a\in\mathbb R^k\}\) be a finite invariant
ansatz.  Suppose all derivatives, contractions, and nonlinear operations
appearing in \(\cE_D\) send \(\cA_{\mathrm{inv}}\) into a finite residual
space with linearly independent basis \(\{e_j\}_{j=1}^{m}\).  Then
\[
 \cE_D(u(a))=\sum_{j=1}^{m}F_j(a)e_j.
\]

\begin{proposition}[Exact scope of an invariant coefficient system]
\label{prop:fcc}
Under the stated closure and linear-independence hypotheses,
\[
 \cE_D(u(a))=0
 \quad\Longleftrightarrow\quad
 F_1(a)=\cdots=F_m(a)=0
\]
for configurations \(u(a)\) inside the chosen ansatz.
\end{proposition}

\begin{proof}
The forward implication follows by taking coefficients.  The converse
follows because vanishing of every coefficient makes the residual zero in
the finite residual space.
\end{proof}

Proposition~\ref{prop:fcc} does not say that every solution of the full PDE is
left-invariant, that time-dependent perturbations remain in the ansatz, or
that the ansatz gives a consistent truncation.  Those are separate
statements.  Nomizu's theorem computes de Rham cohomology from invariant
forms on compact nilmanifolds \cite{Nomizu1954}; it does not imply that all
metrics, connections, nonlinear field configurations, or PDE solutions are
invariant.

Nor does a finite algebraic system automatically select a unique vacuum.
It may have no solutions, several isolated solutions, positive-dimensional
branches, or singular components.  An invertible Jacobian at one solution
gives a locally unique branch under specified parameter variation; it does
not select the integer data, fix an overall scale, or prove global
uniqueness.

Accordingly, ``FCC'' can be retained as a useful name for the residual
coefficient system of a declared invariant ansatz.  It must not be described
as equivalent to the unrestricted higher-dimensional theory without an
independent completeness or consistent-truncation theorem.

\section{Four-dimensional quantities}
\label{sec:outputs}

\subsection{Planck normalization}

For an unwarped product with fixed internal metric and Einstein--Hilbert
term
\[
 S_D^{\mathrm{EH}}
 =
 \frac{1}{2\kappa_D^2}
 \int_{M_4\times X}\sqrt{-g_D}\,R_D,
\]
integration over \(X\) gives
\[
 \frac{1}{\kappa_4^2}
 =
 \frac{\Vol(X,h)}{\kappa_D^2}
\]
before any further convention-dependent rescaling.  Warping, a varying
dilaton, or moduli replace the ordinary volume by a weighted integral and
require an explicit Weyl transformation to four-dimensional Einstein frame.
Thus the formula computes \(\kappa_4\) from supplied
\((\kappa_D,h)\); it does not select either input.

\subsection{Gauge fields and couplings}

If \(K_a\) are Killing vector fields of the internal metric, off-diagonal
metric components along them can produce lower-dimensional vector fields.
The Lie brackets
\[
 [K_a,K_b]=f_{ab}{}^{c}K_c
\]
give the candidate gauge algebra.  In a simple unwarped convention, the
kinetic matrix contains the internal Gram integral
\[
 \cG_{ab}
 \propto
 \frac{1}{\kappa_D^2}
 \int_X
 h(K_a,K_b)\,\dd\mathrm{vol}_h .
\]
The proportionality factor depends on the metric ansatz, generator
normalization, radius factors, and Weyl frame.  Higher-dimensional gauge
connections provide additional gauge sectors by a different mechanism.

An internal isometry therefore makes a gauge interpretation possible, but
it does not select the observed gauge group or its normalized couplings.
Moreover, retaining all isometry vectors need not define a nonlinearly
consistent truncation for an arbitrary compact space
\cite{CveticLuPope2000}.

\subsection{Overlap couplings}

For normalized internal profiles, effective cubic coefficients have the
generic form
\[
 g_{ijk}^{(4)}
 =
 g_D
 \int_X
 \cI\!\left(u_i,u_j,u_k;
 h,\nabla,F,\ldots\right)
 \dd\mathrm{vol}_h ,
\]
where \(\cI\) includes the contractions and background insertions dictated
by the higher-dimensional action.  Selection rules can follow from
representations, parity, index theory, or bundle cohomology.  Numerical
Yukawa or threshold values require the actual normalized wavefunctions,
connections, scales, and renormalization transport.  The mere existence of
an overlap formula is not a value-source theorem.

\section{Relation to current MTT geometry}
\label{sec:status}

\subsection{What is already available}

Current MTT work provides several pieces that are relevant to a future
Kaluza--Klein source map:
\begin{enumerate}[leftmargin=2em]
\item a dimension-neutral Hilbert-bundle architecture with explicit
      operator domains, coherent spectral projectors, and separate
      stabilization and truncation gates;
\item the canonical \(4+6\) physical specialization;
\item the local six-dimensional strain decomposition
      \(\operatorname{Sym}(3)=1\oplus2\oplus3\);
\item a selected finite \(q=79\) carrier with the same rank filtration; and
\item one universal finite differential line with connection- and
      holonomy-preserving pullbacks across the proven finite sectors
      \cite{NeroFoundation}.
\end{enumerate}
These are nontrivial compatibility results.  They make the desired global
map more structured than a guess.

\subsection{What is not yet available}

Rank agreement does not identify the local strain carrier with the vertical
tangent or field bundle of a physical compactification.  The current open
step is a same-source continuum intertwiner that preserves:
\[
\text{bundle transitions},\quad
\text{connection},\quad
\nabla,\quad
\text{measure},\quad
\text{operator domains},\quad
\text{Hessian}.
\]
The physical \(q=79\) program additionally needs the selected visible and
hidden holomorphic bundles and their Hermitian--Yang--Mills connections
before those vertical operators can be executed.  Until then, the finite
rank carrier is not a Kaluza--Klein compactification record.

\subsection{Calabi--Yau, Fu--Yau, and M-theory branches}

The corrected Calabi--Yau paper treats strict Calabi--Yau compactification as
a conditional realization map and separates fixed-background spectral
facts from vacuum selection \cite{NeroCalabiYau}.  The selected \(q=79\)
Fu--Yau program is instead a torsional non-K\"ahler heterotic branch; it
cannot be substituted for a strict Calabi--Yau result merely because both
use six internal dimensions.

The corrected M-theory paper treats the eleven-dimensional low-energy
record and its massless type-IIA circle reduction
\cite{NeroMTheory}.  That circle is not automatically the finite shared MTT
phase line.  This paper supplies the general Kaluza--Klein reduction
interface; it does not re-own either the Calabi--Yau spectral theorem or the
M-theory circle dictionary.

\subsection{Lens--Nil and recursive topology}

Circle, Lens, and Nil language may organize filtrations, bundle towers, or
effective local models.  Literal
\(S^1\times\mathrm{Lens}\times\mathrm{Nil}\) and literal manifold nesting are
not current proof sources for the physical \(q=79\) compactification.  In
particular, the old Lens--Nil and Fu--Yau models have different global
topological data and must not be identified.  A recursive or
circle-fibered description becomes physical only when the relevant total
space, transition functions, connection, and operator compatibility are
specified.

\section{Claim disposition}
\label{sec:audit}

\renewcommand{\arraystretch}{1.2}
\begin{tabularx}{\textwidth}{@{}p{0.27\textwidth}p{0.18\textwidth}X@{}}
\toprule
Version 1 claim & Version 2 status & Resolution \\
\midrule
MTT derives Kaluza--Klein compactification from first principles
& Withdrawn
& MTT currently supplies a conditional source-and-factorization contract. \\

MTT selects \(B_{\rm KK}\), \(R_{\rm KK}\), and the zero-mode spectrum
& Open
& These must be emitted as geometry, scale, and operator data by one
  selected source. \\

\(9\to3\) projection independently proves \(10\to4\) compactification
& Narrowed
& They are spatial and spacetime descriptions of the same supplied
  time-compatible \(4+6\) fibration. \\

The finite FCC is equivalent to the higher-dimensional PDE
& Corrected
& It is equivalent only to the residual equations restricted to a closed
  finite ansatz. \\

Invariant cohomology makes the ansatz complete
& Withdrawn
& Cohomological completeness does not imply field-configuration or PDE
  completeness. \\

A spectral gap justifies zero-mode truncation
& Corrected
& A gap helps invert the omitted linear block; residual and nonlinear bounds
  are also required. \\

Planck, gauge, and KK masses are MTT predictions
& Conditional formulas
& They are computed from supplied action normalizations, geometry, operators,
  and radii. \\

Internal isometries select the observed gauge group
& Open
& Isometries provide candidate gauge vectors; selection and nonlinear
  consistency require additional data. \\

Iwasawa and Lens--Nil examples prove physical compactification selection
& Retired as proof sources
& They may remain auxiliary invariant models but do not establish the
  current physical \(q=79\) branch. \\

MTT gap data fix families, thresholds, and observed couplings
& Removed
& Those claims require selected normalized modes, source values, and
  renormalization transport. \\
\bottomrule
\end{tabularx}

\section{Completion contract}
\label{sec:completion}

A selected MTT Kaluza--Klein realization will be established only when one
source supplies and one independent verification checks the following:
\begin{enumerate}[leftmargin=2em]
\item a global \(D\)-dimensional Lorentzian fibration with compact internal
      geometry, transition maps, horizontal connection, spin or
      spin\(^{c}\) data, and all required bundles;
\item the higher-dimensional field content, action, symmetries, sources,
      normalizations, and declared classical or quantum order;
\item the vertical operators with domains, measures, adjoints, boundary
      conditions, normalized modes, and scale conventions;
\item the selected background and a proof that it solves the supplied
      equations, rather than only a subset of projected coefficients;
\item an exact nonlinear uplift or a residual certificate of the form in
      Proposition~\ref{prop:residual}, including propagation to the stated
      observables;
\item a source-preserving map from the MTT strain and shared-line data to the
      physical vertical bundles, connections, derivatives, and Hessians;
\item normalized four-dimensional Planck, gauge, mass, chirality, and
      interaction data derived from that same source; and
\item a clear separation between construction inputs, fitted values,
      profile checks, and held-out predictions.
\end{enumerate}

The first five items would establish a controlled Kaluza--Klein reduction.
The sixth would make it an MTT-sourced reduction.  The final two would be
needed before phenomenological equivalence or prediction could be claimed.

\section{Discussion}
\label{sec:discussion}

\subsection{Why this correction strengthens the program}

The corrected paper asks less rhetorically and more mathematically.  It no
longer treats the existence of familiar lower formulas as evidence that
their inputs were derived.  Instead, it turns the missing source into a
finite list of typed obligations.  This is useful even before closure:
different candidate compactifications can be compared row by row, and a
failure in geometry, action, truncation, or normalization cannot be hidden
inside a generic projection symbol.

The distinction between spectrum and truncation is especially important.
Compactness gives a tower.  It does not say that the light part of the tower
is closed under interactions.  The residual certificate identifies the
quantities that a computation must actually bound.  It also gives a natural
place for the MTT fixed-point and semigroup machinery: those methods may
control \(\gamma_N\), \(L_\rho\), and \(r_N\), but only after the physical
operator and its source have been identified.

\subsection{What could become distinctive}

Traditional Kaluza--Klein theory begins after a higher-dimensional record is
chosen.  MTT asks one level upstream whether the same upper object sources
the internal geometry, shared line, operators, action, and admissible low
sector.  A connection-preserving intertwiner and upper-action map would not
replace Kaluza--Klein theory; they would explain why a particular record is
the correct lower image and why its truncation is admissible.

That is the credible advantage of the MTT route.  Component counts and
replayed four-dimensional values remain insufficient: the upstream source
must reduce independent choices, not rename them.

\section{Conclusion}

Kaluza--Klein reduction is exact mathematics once its higher-dimensional
record and consistency hypotheses are fixed.  On a closed compact fiber,
self-adjoint elliptic operators give discrete towers, and the tower masses
are internal eigenvalues with the appropriate metric scaling.  Internal
isometries and normalized overlap integrals can generate gauge fields and
effective couplings.  Finite invariant ansatzes can turn their restricted
equations into exact algebraic systems.

None of those facts selects the internal space, radius, action, gauge group,
or nonlinear truncation.  Version 2 therefore replaces the former
first-principles claim by a conditional MTT realization theorem and an
explicit truncation certificate.  It reconciles the spatial \(9\to3\)
language with the canonical \(10\to4\) fibration, keeps the shared circle as
common phase/holonomy data at its proven level, and separates local
\(1+2+3\) rank agreement from the still-open physical continuum
intertwiner.

The route forward is sharp: select one physical vertical geometry and
action, preserve their connection and Hessian under the MTT source map,
execute the appropriate internal operators, and certify the retained sector
by exact consistency or quantitative residual bounds.  That would turn this
contract into a selected MTT Kaluza--Klein realization.  Until then it proves
compatibility and a rigorous proof interface, not extra-dimensional
phenomenology.

\bibliographystyle{unsrt}

\enlargethispage{3\baselineskip}

% BEGIN MTT MANAGED COMPUTATIONAL EVIDENCE
\section*{Computational Evidence and Reproducibility}
The exact q79 arithmetic, finite rank-two Cech witness, and rank-two Wiener-contraction certificate are contextual evidence for an adjacent selected heterotic branch. They do not directly construct the physical compact six-manifold, higher-dimensional action, nonlinear uplift, normalized vertical mode system, Kaluza-Klein radius, gauge group, consistent truncation, or the connection- and Hessian-preserving MTT-to-Kaluza-Klein realization map.

The referenced rows are frozen to the curated results repository state identified below. Hashes are grouped in eight-character blocks for line breaking.
\begin{quote}\small
\textbf{Repository:} \url{https://github.com/PeterNero/mtt-results-repro}\\
\textbf{Commit:} \texttt{31247ebb 5c22f3fb b5443024 365433c6 ee0bff4a}\\
\textbf{Manifest:} \path{release/result_manifest.json}\\
\textbf{Manifest SHA-256:}\\
\texttt{fb399689 60b00584 631dbf53 1a708e18}\\
\texttt{ef928d6b 6d935119 c185d7f6 32b1e7cd}
\end{quote}

Tier labels are quoted verbatim from that manifest. A row used directly supports only the specific computational statement identified above; a corpus-state cross-check does not prove this paper's local theorems; and an open row is evidence of an unresolved obligation, never of closure.

\par\begingroup\dimen0=\pagegoal\advance\dimen0 by -\pagetotal\ifdim\dimen0<7\baselineskip\newpage\fi\endgroup
\paragraph{Corpus-state cross-checks.}
\begin{itemize}
\setlength{\itemsep}{0.25em}
\item \path{A19/hym_wiener_contraction} (\emph{numeric certified}).\par Weighted-theta Fourier-tail and Wiener contraction certificate.
\item \path{A07/literal_cech_witness} (\emph{derived exact}).\par Literal 81-entry, 729-cocycle finite Cech witness.
\item \path{A11/q79_exact_audit} (\emph{derived exact}).\par Executable q=79 exact-branch audit.
\item \path{A11/q79_exact_theorem} (\emph{derived exact}).\par CRT q=79 theorem on the selected exact branch.
\end{itemize}

No imported row changes theorem ownership or promotes a neighboring claim: all local statements retain their stated hypotheses, domains, and limitations.
% END MTT MANAGED COMPUTATIONAL EVIDENCE

\bibliography{main}

\end{document}
